\subsubsection{Diameter and average path length}

The same Python package \texttt{NetworkX} has many useful functions. Once the Watts-Strogatz graph has been generated using the \texttt{nx.watts\_strogatz\_graph} function, the following functions are then used: \texttt{nx.diameter()} and \texttt{nx.average\_shortest\_path\_length()}. As the names of the functions indicate, the \texttt{nx.diameter()} returns the diameter (as defined in equation \ref{eq:WSdiameter}) and the latter function computes the average path length (as defined in equation \ref{eq:WSavgshortestpath}). In the case of an disconnected graph, the diameter and average path length would be computed for the largest connected component of the graph.

The computations are repeated 1000 times for each value of $p$, with a new seed each times. For every iteration, a new graph is generated and the relevant values are saved in a \texttt{NumPy}-array. This will then be used to produce the results visible in section \ref{sec:WS_Diam&ASPL_result}.
